\scanpage{27}
% archive foliation: 26
\uline{Corollaire 1} Dans $\mathrm{Hot}(\underline W)$, les sommes directes finies existent, et les foncteurs canoniques $(\mathrm{Cat})\longrightarrow\mathrm{Hot}(\underline W)$, $\widehat A\longrightarrow\mathrm{Hot}_A,\longrightarrow\mathrm{Hot}(\underline W)$ y commutent.

\margin{\uline{finies} inutile, \unsure{par argument plus bas}}
\margin{pas de conditions sur $A$ ici! \diagram{\begin{tikzcd}&{}\\{}\arrow[ur]&\end{tikzcd}}}

On est ramené au cas de $(\mathrm{Cat})$ (vu que $i_A:\widehat A\longrightarrow(\mathrm{Cat})$ commute aux $\varinjlim$ \unsure{donc aux} sommes), via le

\margin{$|$}\uline{Lemme} Soit $(\underline M,W)$ avec $W\subset\mathrm{Fl}(\underline M)$, $W$ permettant calcul des fractions à droite. \deleted{Si dans $M$ les} \inserted{Alors} \inserted{le} foncteur canonique $M\longrightarrow W^{-1}M$ commute aux $\varinjlim$ finies, \deleted{\unsure{les}} et en particulier aux sommes finies.

On appliquera ceci à $(\overline{\mathrm{Cat}},\overline{\underline W})$, il faut voir que les sommes finies y existent (sinon le lemme n'apprend pas \unsure{gd} chose!), et que $\mathrm{Cat}\longrightarrow\overline{\mathrm{Cat}}$ y commute. \deleted{Donc} Mais on a de façon générale
